% ============================================================
% Question Bank: ch05-04 Solving Multi-Step Problems with Rational Numbers
% Topic Title: Solving Multi-Step Problems with Rational Numbers
% CCSS: 7.EE.B.3
% Questions: 30 (18 MC + 7 SA + 5 Graphical)
% ============================================================

% --- Multiple Choice Questions (18) ---

\begin{practiceQuestion}{5-4-q01}{mc}
\begin{questionText}
Solve $0.5x + 3 = 8$.
\end{questionText}
\choiceA{$x = 10$}
\choiceB{$x = 22$}
\choiceC{$x = 5.5$}
\choiceD{$x = 2.5$}
\correctAnswer{A}
\explanation{Subtract $3$: $0.5x = 5$. Divide by $0.5$: $x = 10$. Check: $0.5(10) + 3 = 5 + 3 = 8$ \checkmark.}
\end{practiceQuestion}

\begin{practiceQuestion}{5-4-q02}{mc}
\begin{questionText}
Solve $\dfrac{x}{3} + 5 = 9$.
\end{questionText}
\choiceA{$x = 42$}
\choiceB{$x = 12$}
\choiceC{$x = \dfrac{4}{3}$}
\choiceD{$x = 1\dfrac{1}{3}$}
\correctAnswer{B}
\explanation{Subtract $5$: $\frac{x}{3} = 4$. Multiply by $3$: $x = 12$. Check: $\frac{12}{3} + 5 = 4 + 5 = 9$ \checkmark.}
\end{practiceQuestion}

\begin{practiceQuestion}{5-4-q03}{mc}
\begin{questionText}
Solve $0.25n - 1.5 = 3.5$.
\end{questionText}
\choiceA{$n = 8$}
\choiceB{$n = 20$}
\choiceC{$n = 14$}
\choiceD{$n = 2$}
\correctAnswer{B}
\explanation{Add $1.5$: $0.25n = 5$. Divide by $0.25$: $n = 20$. Check: $0.25(20) - 1.5 = 5 - 1.5 = 3.5$ \checkmark.}
\end{practiceQuestion}

\begin{practiceQuestion}{5-4-q04}{mc}
\begin{questionText}
Which is the best first step to solve $\dfrac{x}{4} + \dfrac{1}{2} = 3$?
\end{questionText}
\choiceA{Subtract $\frac{1}{2}$ from both sides}
\choiceB{Multiply every term by $4$}
\choiceC{Divide both sides by $4$}
\choiceD{Add $\frac{1}{2}$ to both sides}
\correctAnswer{B}
\explanation{Multiplying every term by $4$ (the LCD) clears the fractions: $x + 2 = 12$.}
\end{practiceQuestion}

\begin{practiceQuestion}{5-4-q05}{mc}
\begin{questionText}
Solve $\dfrac{2x}{5} - 1 = 3$.
\end{questionText}
\choiceA{$x = 5$}
\choiceB{$x = 10$}
\choiceC{$x = 4$}
\choiceD{$x = 20$}
\correctAnswer{B}
\explanation{Add $1$: $\frac{2x}{5} = 4$. Multiply by $5$: $2x = 20$. Divide by $2$: $x = 10$. Check: $\frac{2(10)}{5} - 1 = 4 - 1 = 3$ \checkmark.}
\end{practiceQuestion}

\begin{practiceQuestion}{5-4-q06}{mc}
\begin{questionText}
Solve $-0.4y + 2 = 6$.
\end{questionText}
\choiceA{$y = 10$}
\choiceB{$y = -10$}
\choiceC{$y = 20$}
\choiceD{$y = -20$}
\correctAnswer{B}
\explanation{Subtract $2$: $-0.4y = 4$. Divide by $-0.4$: $y = -10$. Check: $-0.4(-10) + 2 = 4 + 2 = 6$ \checkmark.}
\end{practiceQuestion}

\begin{practiceQuestion}{5-4-q07}{mc}
\begin{questionText}
Solve $\dfrac{3n}{4} + 6 = 15$.
\end{questionText}
\choiceA{$n = 28$}
\choiceB{$n = 12$}
\choiceC{$n = 7$}
\choiceD{$n = 36$}
\correctAnswer{B}
\explanation{Subtract $6$: $\frac{3n}{4} = 9$. Multiply by $4$: $3n = 36$. Divide by $3$: $n = 12$. Check: $\frac{3(12)}{4} + 6 = 9 + 6 = 15$ \checkmark.}
\end{practiceQuestion}

\begin{practiceQuestion}{5-4-q08}{mc}
\begin{questionText}
A recipe calls for $\dfrac{3}{4}$ cup of sugar per batch. You already used $\dfrac{1}{2}$ cup. You have $2\dfrac{3}{4}$ cups total. The equation $\dfrac{3}{4}b + \dfrac{1}{2} = 2\dfrac{3}{4}$ finds how many more batches $b$ you can make. What is $b$?
\end{questionText}
\choiceA{$b = 3$}
\choiceB{$b = 4$}
\choiceC{$b = 2$}
\choiceD{$b = 3\dfrac{2}{3}$}
\correctAnswer{A}
\explanation{Subtract $\frac{1}{2}$: $\frac{3}{4}b = 2\frac{1}{4} = \frac{9}{4}$. Multiply by $\frac{4}{3}$: $b = \frac{9}{4} \times \frac{4}{3} = 3$.}
\end{practiceQuestion}

\begin{practiceQuestion}{5-4-q09}{mc}
\begin{questionText}
Solve $1.2x + 0.8 = 4.4$.
\end{questionText}
\choiceA{$x = 3$}
\choiceB{$x = 4.33$}
\choiceC{$x = 3.6$}
\choiceD{$x = 2.5$}
\correctAnswer{A}
\explanation{Subtract $0.8$: $1.2x = 3.6$. Divide by $1.2$: $x = 3$. Check: $1.2(3) + 0.8 = 3.6 + 0.8 = 4.4$ \checkmark.}
\end{practiceQuestion}

\begin{practiceQuestion}{5-4-q10}{mc}
\begin{questionText}
Solve $\dfrac{x}{6} - \dfrac{2}{3} = \dfrac{1}{2}$.
\end{questionText}
\choiceA{$x = 3$}
\choiceB{$x = 7$}
\choiceC{$x = -1$}
\choiceD{$x = 10$}
\correctAnswer{B}
\explanation{Multiply every term by $6$: $x - 4 = 3$. Add $4$: $x = 7$. Check: $\frac{7}{6} - \frac{2}{3} = \frac{7}{6} - \frac{4}{6} = \frac{3}{6} = \frac{1}{2}$ \checkmark.}
\end{practiceQuestion}

\begin{practiceQuestion}{5-4-q11}{mc}
\begin{questionText}
Solve $-1.5x + 4.5 = -3$.
\end{questionText}
\choiceA{$x = -5$}
\choiceB{$x = 1$}
\choiceC{$x = 5$}
\choiceD{$x = -1$}
\correctAnswer{C}
\explanation{Subtract $4.5$: $-1.5x = -7.5$. Divide by $-1.5$: $x = 5$. Check: $-1.5(5) + 4.5 = -7.5 + 4.5 = -3$ \checkmark.}
\end{practiceQuestion}

\begin{practiceQuestion}{5-4-q12}{mc}
\begin{questionText}
Solve $\dfrac{n}{2} + \dfrac{n}{3} = 10$.
\end{questionText}
\choiceA{$n = 15$}
\choiceB{$n = 12$}
\choiceC{$n = 6$}
\choiceD{$n = 20$}
\correctAnswer{B}
\explanation{Multiply every term by $6$: $3n + 2n = 60$. Combine: $5n = 60$. Divide by $5$: $n = 12$. Check: $\frac{12}{2} + \frac{12}{3} = 6 + 4 = 10$ \checkmark.}
\end{practiceQuestion}

\begin{practiceQuestion}{5-4-q13}{mc}
\begin{questionText}
A shirt is on sale for $\dfrac{2}{3}$ of its original price. After using a $\$5$ coupon, you pay $\$15$. What equation finds the original price $p$?
\end{questionText}
\choiceA{$\dfrac{2}{3}p + 5 = 15$}
\choiceB{$p - \dfrac{2}{3} - 5 = 15$}
\choiceC{$\dfrac{2}{3}p - 5 = 15$}
\choiceD{$\dfrac{2}{3}(p - 5) = 15$}
\correctAnswer{C}
\explanation{Sale price is $\frac{2}{3}p$. After the $\$5$ coupon: $\frac{2}{3}p - 5 = 15$.}
\end{practiceQuestion}

\begin{practiceQuestion}{5-4-q14}{mc}
\begin{questionText}
Solve $0.6x - 2.4 = 1.2$.
\end{questionText}
\choiceA{$x = 6$}
\choiceB{$x = -2$}
\choiceC{$x = 2$}
\choiceD{$x = 0.6$}
\correctAnswer{A}
\explanation{Add $2.4$: $0.6x = 3.6$. Divide by $0.6$: $x = 6$. Check: $0.6(6) - 2.4 = 3.6 - 2.4 = 1.2$ \checkmark.}
\end{practiceQuestion}

\begin{practiceQuestion}{5-4-q15}{mc}
\begin{questionText}
Solve $\dfrac{5x}{2} + 3 = 18$.
\end{questionText}
\choiceA{$x = 3$}
\choiceB{$x = 6$}
\choiceC{$x = 42$}
\choiceD{$x = 15$}
\correctAnswer{B}
\explanation{Subtract $3$: $\frac{5x}{2} = 15$. Multiply by $2$: $5x = 30$. Divide by $5$: $x = 6$. Check: $\frac{5(6)}{2} + 3 = 15 + 3 = 18$ \checkmark.}
\end{practiceQuestion}

\begin{practiceQuestion}{5-4-q16}{mc}
\begin{questionText}
A student estimates that the solution to $0.48x + 3.1 = 12.7$ is about $x = 20$. Is this estimate reasonable?
\end{questionText}
\choiceA{Yes, the exact answer is $20$}
\choiceB{Yes, the exact answer is close to $20$}
\choiceC{No, the exact answer is about $10$}
\choiceD{No, the exact answer is about $33$}
\correctAnswer{A}
\explanation{Subtract $3.1$: $0.48x = 9.6$. Divide by $0.48$: $x = 20$. The estimate is exactly correct.}
\end{practiceQuestion}

\begin{practiceQuestion}{5-4-q17}{mc}
\begin{questionText}
Solve $\dfrac{x}{5} + \dfrac{3}{10} = \dfrac{7}{10}$.
\end{questionText}
\choiceA{$x = 4$}
\choiceB{$x = 5$}
\choiceC{$x = 1$}
\choiceD{$x = 2$}
\correctAnswer{D}
\explanation{Multiply every term by $10$: $2x + 3 = 7$. Subtract $3$: $2x = 4$. Divide by $2$: $x = 2$. Check: $\frac{2}{5} + \frac{3}{10} = \frac{4}{10} + \frac{3}{10} = \frac{7}{10}$ \checkmark.}
\end{practiceQuestion}

\begin{practiceQuestion}{5-4-q18}{mc}
\begin{questionText}
Solve $-0.2n + 5.6 = 4$.
\end{questionText}
\choiceA{$n = 8$}
\choiceB{$n = -8$}
\choiceC{$n = 48$}
\choiceD{$n = -48$}
\correctAnswer{A}
\explanation{Subtract $5.6$: $-0.2n = -1.6$. Divide by $-0.2$: $n = 8$. Check: $-0.2(8) + 5.6 = -1.6 + 5.6 = 4$ \checkmark.}
\end{practiceQuestion}

% --- Short Answer Questions (7) ---

\begin{practiceQuestion}{5-4-q19}{sa}
\begin{questionText}
Solve $\dfrac{3x}{4} - 2 = 7$.
\end{questionText}
\correctAnswer{$x = 12$}
\explanation{Add $2$: $\frac{3x}{4} = 9$. Multiply by $4$: $3x = 36$. Divide by $3$: $x = 12$. Check: $\frac{3(12)}{4} - 2 = 9 - 2 = 7$ \checkmark.}
\end{practiceQuestion}

\begin{practiceQuestion}{5-4-q20}{sa}
\begin{questionText}
Solve $0.75y + 2.25 = 8.25$.
\end{questionText}
\correctAnswer{$y = 8$}
\explanation{Subtract $2.25$: $0.75y = 6$. Divide by $0.75$: $y = 8$. Check: $0.75(8) + 2.25 = 6 + 2.25 = 8.25$ \checkmark.}
\end{practiceQuestion}

\begin{practiceQuestion}{5-4-q21}{sa}
\begin{questionText}
Solve $\dfrac{x}{2} + \dfrac{x}{4} = 9$.
\end{questionText}
\correctAnswer{$x = 12$}
\explanation{Multiply every term by $4$: $2x + x = 36$. Combine: $3x = 36$. Divide by $3$: $x = 12$. Check: $\frac{12}{2} + \frac{12}{4} = 6 + 3 = 9$ \checkmark.}
\end{practiceQuestion}

\begin{practiceQuestion}{5-4-q22}{sa}
\begin{questionText}
Solve $-0.3x + 4.2 = 1.8$.
\end{questionText}
\correctAnswer{$x = 8$}
\explanation{Subtract $4.2$: $-0.3x = -2.4$. Divide by $-0.3$: $x = 8$. Check: $-0.3(8) + 4.2 = -2.4 + 4.2 = 1.8$ \checkmark.}
\end{practiceQuestion}

\begin{practiceQuestion}{5-4-q23}{sa}
\begin{questionText}
A store marks down a jacket by $\dfrac{1}{4}$ of its original price $p$, then subtracts an additional $\$8$. The sale price is $\$37$. Find the original price.
\end{questionText}
\correctAnswer{$p = \$60$}
\explanation{After $\frac{1}{4}$ markdown: $p - \frac{p}{4} = \frac{3p}{4}$. After subtracting $\$8$: $\frac{3p}{4} - 8 = 37$. Add $8$: $\frac{3p}{4} = 45$. Multiply by $\frac{4}{3}$: $p = 60$.}
\end{practiceQuestion}

\begin{practiceQuestion}{5-4-q24}{sa}
\begin{questionText}
Solve $\dfrac{5n}{6} + \dfrac{1}{3} = 3$.
\end{questionText}
\correctAnswer{$n = \dfrac{16}{5}$ or $3\dfrac{1}{5}$}
\explanation{Multiply every term by $6$: $5n + 2 = 18$. Subtract $2$: $5n = 16$. Divide by $5$: $n = \frac{16}{5} = 3\frac{1}{5}$.}
\end{practiceQuestion}

\begin{practiceQuestion}{5-4-q25}{sa}
\begin{questionText}
Solve $2.5x - 7.5 = 10$.
\end{questionText}
\correctAnswer{$x = 7$}
\explanation{Add $7.5$: $2.5x = 17.5$. Divide by $2.5$: $x = 7$. Check: $2.5(7) - 7.5 = 17.5 - 7.5 = 10$ \checkmark.}
\end{practiceQuestion}

% --- Graphical Questions (3 GMC + 2 GSA) ---

\begin{practiceQuestion}{5-4-q26}{gmc}
\begin{questionText}
The table shows input and output values for a function. What is the rule?

\begin{center}
\begin{tabular}{|c|c|}
\hline
\textbf{Input ($x$)} & \textbf{Output ($y$)} \\
\hline
$2$ & $2.5$ \\
\hline
$4$ & $3.5$ \\
\hline
$6$ & $4.5$ \\
\hline
$10$ & $6.5$ \\
\hline
\end{tabular}
\end{center}
\end{questionText}
\choiceA{$y = 0.5x + 1.5$}
\choiceB{$y = x - 0.5$}
\choiceC{$y = 0.25x + 2$}
\choiceD{$y = 2x - 1.5$}
\correctAnswer{A}
\explanation{From $x = 2$ to $x = 4$, $y$ increases by $1$ while $x$ increases by $2$, so rate $= 0.5$. When $x = 2$: $0.5(2) + 1.5 = 2.5$ \checkmark.}
\end{practiceQuestion}

\begin{practiceQuestion}{5-4-q27}{gmc}
\begin{questionText}
Use the fraction bar model to solve: $\dfrac{2}{3}x = 8$. What is $x$?

\begin{center}
\begin{tikzpicture}[scale=0.6]
  % Full bar = x
  \draw[thick,fill=funBlue!15] (0,2) rectangle (12,3);
  \node at (6,2.5) {$x$};
  % 2/3 of bar highlighted
  \draw[thick,fill=funOrange!30] (0,0) rectangle (8,1);
  \node at (4,0.5) {$\frac{2}{3}x = 8$};
  \draw[thick,fill=gray!15] (8,0) rectangle (12,1);
  \node at (10,0.5) {$?$};
  % Dividers
  \draw[dashed] (4,0) -- (4,1);
  \draw[dashed] (4,2) -- (4,3);
  \draw[dashed] (8,0) -- (8,3);
  \node[below] at (2,-0.2) {$\frac{1}{3}$};
  \node[below] at (6,-0.2) {$\frac{1}{3}$};
  \node[below] at (10,-0.2) {$\frac{1}{3}$};
\end{tikzpicture}
\end{center}
\end{questionText}
\choiceA{$x = 12$}
\choiceB{$x = 16$}
\choiceC{$x = 10$}
\choiceD{$x = 5\dfrac{1}{3}$}
\correctAnswer{A}
\explanation{If $\frac{2}{3}x = 8$, then each $\frac{1}{3}$ of $x$ is $4$. So $x = 3 \times 4 = 12$.}
\end{practiceQuestion}

\begin{practiceQuestion}{5-4-q28}{gmc}
\begin{questionText}
The thermometer shows the current temperature. The temperature was $-3.5\degree$F this morning and rose $t$ degrees each hour for $4$ hours to reach the current reading. What is $t$?

\begin{center}
\begin{tikzpicture}[scale=0.5]
  % Thermometer
  \draw[thick] (0,0) -- (0,10) arc(180:0:0.5) -- (1,0) arc(0:-180:0.5);
  \fill[red!60] (0.1,0) rectangle (0.9,6.5);
  % Scale
  \foreach \y/\val in {0/-5, 2/0, 4/5, 6/10, 8/15, 10/20} {
    \draw (-0.3,\y) -- (0,\y);
    \node[left] at (-0.4,\y) {$\val\degree$F};
  }
  % Arrow to current
  \draw[thick,->,funBlue] (1.8,6.5) -- (1.1,6.5);
  \node[right] at (1.8,6.5) {Current: $8.5\degree$F};
\end{tikzpicture}
\end{center}
\end{questionText}
\choiceA{$t = 3$}
\choiceB{$t = 5$}
\choiceC{$t = 1.25$}
\choiceD{$t = 2$}
\correctAnswer{A}
\explanation{Equation: $-3.5 + 4t = 8.5$. Add $3.5$: $4t = 12$. Divide by $4$: $t = 3$.}
\end{practiceQuestion}

\begin{practiceQuestion}{5-4-q29}{gsa}
\begin{questionText}
The bar graph shows how much three students earned. Together they earned a total of $\$31.50$. Student C earned $\dfrac{1}{2}$ as much as Student A. Find how much Student C earned.

\begin{center}
\begin{tikzpicture}[scale=0.5]
  \draw[thick,->] (0,0) -- (0,11) node[above] {Dollars (\$)};
  \draw[thick,->] (0,0) -- (10,0);
  % Bars
  \draw[thick,fill=funBlue!40] (1,0) rectangle (3,8);
  \node[below] at (2,-0.3) {A};
  \node at (2,8.5) {$\$12$};
  \draw[thick,fill=funOrange!40] (4,0) rectangle (6,0);
  \node[below] at (5,-0.3) {B};
  \node at (5,1) {?};
  \draw[thick,fill=funGreen!40] (7,0) rectangle (9,0);
  \node[below] at (8,-0.3) {C};
  \node at (8,1) {?};
  % Y-axis labels
  \foreach \y in {0,2,...,10} {
    \draw (-0.2,\y) -- (0,\y);
    \node[left] at (-0.3,\y) {$\y$};
  }
\end{tikzpicture}
\end{center}

Student A earned $\$12$, and Student C earned $\dfrac{1}{2}$ of Student A's earnings.
\end{questionText}
\correctAnswer{$\$6$}
\explanation{Student C $= \frac{1}{2} \times 12 = \$6$. Check: Student B $= 31.50 - 12 - 6 = \$13.50$. Total: $12 + 13.50 + 6 = 31.50$ \checkmark.}
\end{practiceQuestion}

\begin{practiceQuestion}{5-4-q30}{gsa}
\begin{questionText}
A number line shows a point at $x$. The equation $\dfrac{x}{4} + 1.5 = 5$ tells you how to find $x$. Solve and identify where $x$ is on the number line.

\begin{center}
\begin{tikzpicture}[scale=0.45]
  \draw[thick,->] (-2,0) -- (18,0);
  \foreach \x in {0,2,...,16} {
    \draw (\x,0.2) -- (\x,-0.2) node[below] {$\x$};
  }
  \node[above,funBlue] at (14,0.5) {$x = \;?$};
  \draw[thick,->,funBlue] (14,0.4) -- (14,0.1);
\end{tikzpicture}
\end{center}
\end{questionText}
\correctAnswer{$x = 14$}
\explanation{Subtract $1.5$: $\frac{x}{4} = 3.5$. Multiply by $4$: $x = 14$. The point is at $14$ on the number line.}
\end{practiceQuestion}
